\SourceLeaf{102}%
schehen. Man kann, wie es die Berechner bisher gemacht haben, den sogenannten wahrscheinlichen (nicht wahrschein\textit{lichsten}) Fehler zum Maassstabe wählen, welches derjenige ist, über welchen hinaus alle möglichen Fehler zusammen noch eben so viele Wahrscheinlichkeit haben, wie alle diesseits liegenden zusammen; allein es wird \textit{weit vortheilhafter} sein, zu diesem Zweck statt des wahrscheinlichen Fehlers den \textit{mittlern} zu gebrauchen, vorausgesetzt, dass man diesen an sich noch schwankenden Begriff auf die rechte Art bestimmt. Man lege jedem Fehler ein von seiner Grösse abhängendes Moment bei, multiplicire das Moment jedes möglichen Fehlers in dessen Wahrscheinlichkeit und addire die Producte: der Fehler, dessen Moment diesem Aggregat gleich ist, wird als mittlerer betrachtet werden müssen. Allein welche Function der Grösse des Fehlers wir für dessen Moment wählen wollen, bleibt wieder unsrer \textit{Willkür} überlassen, wenn nur der Werth derselben immer positiv ist, und für grössere Fehler grösser als für kleinere. Der Verf. hat die einfachste Function dieser Art gewählt, nemlich das Quadrat; diese Wahl ist aber noch mit manchen andern höchst wesentlichen Vortheilen verknüpft, die bei keiner andern statt finden. Denn sonst könnte auch jede andere Potenz mit geraden Exponenten gebraucht werden, und je grösser dieser Exponent gewählt würde, desto näher würde man dem Princip kommen, wo bloss die äussersten Fehler zum Maassstabe der Genauigkeit dienen. Gegen die Art, wie ein grosser Geometer den Begriff des mittlern Fehlers genommen hat, indem er die Momente der Fehler diesen gleich setzt, wenn sie positiv sind, und die ihnen entgegengesetzten Grössen dafür gebraucht, wenn sie negativ sind, lässt sich bemerken, dass dabei gegen die mathematische Continuität angestossen wird, dass sie so gut wie jede andere auch willkürlich gewählt ist, dass die Resultate viel weniger einfach und genugthuend ausfallen, und dass es auch an sich schon natürlicher scheint, das Moment der Fehler in einem stärkern Verhältniss, wie diese selbst, wachsen zu lassen, indem man sich gewiss lieber den einfachen Fehler zweimal, als den doppelten einmal gefallen lässt.

Diese Erläuterungen mussten vorangeschickt werden, wenn auch nur etwas von dem Inhalt der Untersuchung hier angeführt werden sollte, wovon die gegenwärtige Abhandlung die erste Abtheilung ausmacht.

Wenn die Grössen, deren Werthe durch Beobachtungen gefunden sind, mit einer gleichen Anzahl unbekannter Grössen auf eine bekannte Art zusammenhangen, so lassen sich, allgemein zu reden, die Werthe der unbekannten Grössen
\SourceLeaf{103}
aus den Beobachtungen durch Rechnung ableiten. Freilich werden jene Werthe auch nur näherungsweise richtig sein, in so fern die Beobachtungen es waren: allein die Wahrscheinlichkeitsrechnung hat nichts dabei zu thun, als die Unsicherheit jener Bestimmungen zu würdigen, indem sie die der Beobachtungen voraussetzt. Ist die Anzahl der unbekannten Grössen grösser als die der Beobachtungen, so lassen sich jene aus diesen noch gar nicht bestimmen. Allein wenn die Anzahl der unbekannten Grössen kleiner ist, als die der Beobachtungen, so ist die Aufgabe mehr als bestimmt: es sind dann unendlich viele Combinationen möglich, um aus den Beobachtungen die unbekannten Grössen abzuleiten, die freilich alle zu einerlei Resultaten führen müssten, wenn die Beobachtungen absolute Genauigkeit hätten, aber unter den obwaltenden Umständen mehr oder weniger von einander abweichende Resultate hervorbringen. Aus dieser ins Unendliche gehenden Mannichfaltigkeit von Combinationen die zweckmässigste auszuwählen, d. i. diejenige, wobei die Unsicherheit der Resultate die möglich kleinste wird, ist unstreitig eine der wichtigsten Aufgaben bei der Anwendung der Mathematik auf die Naturwissenschaften.

Der Verfasser gegenwärtiger Abhandlung, welcher im Jahr 1797 diese Aufgabe nach den Grundsätzen der Wahrscheinlichkeitsrechnung zuerst untersuchte, fand bald, dass die Ausmittelung der \textit{wahrscheinlichsten} Werthe der unbekannten Grösse unmöglich sei, wenn nicht die Function, die die Wahrscheinlichkeit der Fehler darstellt, bekannt ist. In so fern sie dies aber nicht ist, bleibt nichts übrig, als hypothetisch eine solche Function anzunehmen. Es schien ihm das natürlichste, zuerst den umgekehrten Weg einzuschlagen und die Function zu suchen, die zum Grunde gelegt werden muss, wenn eine allgemein als gut anerkannte Regel für den einfachsten aller Fälle daraus hervorgehen soll, die nemlich, dass das arithmetische Mittel aus mehreren für eine und dieselbe unbekannte Grösse durch Beobachtungen von gleicher Zuverlässigkeit gefundenen Werthen als der wahrscheinlichste betrachtet werden müsse. Es ergab sich daraus, dass die Wahrscheinlichkeit eines Fehlers \(x\), einer Exponentialgrösse von der Form \(e^{-hhxx}\) proportional angenommen werden müsse, und dass dann gerade diejenige Methode, auf die er schon einige Jahre zuvor durch andere Betrachtungen gekommen war, allgemein nothwendig werde. Diese Methode, welche er nachher besonders seit 1801 bei allerlei astronomischen Rechnungen fast täglich anzuwenden Gelegenheit hatte, und auf welche auch \textsc{Legendre} inzwischen gekommen war,
